%% notes-tex/w019-strain-build-list/sections/1-terms.tex
%% [[experiments.W019-echo-crispr-array.build-list]]
%% https://github.com/Mjvolk3/torchcell/tree/main/notes-tex/w019-strain-build-list/sections/1-terms.tex
%%
%% SOURCE: hand-authored. Definitions only, no data. The column names defined
%% here are produced by build_list_tables.py; the numbers they describe are not
%% repeated in this section.

\section{Terms\secstatus{todo}}\label{sec:terms}

\begin{description}[leftmargin=1.2em,itemsep=2pt,topsep=2pt]

\item[boot SE] The standard error of a strain's mean fitness, from resampling
  the replicate plates that strain was measured on. It is smaller than
  \emph{plate SD} because it is the error on a mean rather than the spread of
  the values that mean was taken over.

\item[fitness] Colony growth relative to the wild type on the same plate, so 1
  is wild-type growth.

\item[plate SD] The standard deviation of a strain's fitness across the
  replicate plates it was measured on.

\item[routes] How many different ways the bench can make a triple from strains
  that already exist. Every triple is built by crossing one existing double with
  the remaining single. \emph{routes} 1 means only one of the triple's three
  pairs is in the freezer, so there is a single way in, and a failed cross
  blocks that triple until the matching new double from
  Table~\ref{tab:new-doubles} is done. \emph{routes} 2 means two of its pairs
  exist, so a failed cross has an immediate second path.

\item[tier] How a double was classed when run 4 was designed.
  \texttt{validation} reproduces a measurement someone has already published,
  \texttt{coverage} was added so that later triples would have a parent to be
  built from, and \texttt{novel} had not been measured by anyone.

\end{description}
